% Preamble
\documentclass[a4paper,11pt]{article}
\title{}
\author{
    André Cardoso 50\%\\
    andremacardoso@ua.pt\\
    108269\\
    \and
    Tiago Figueiredo 50\%\\
    tiago.a.figueiredo@ua.pt\\
    107263\\
}
\date{09 de Janeiro de 2023}

%packages
\usepackage{graphicx}
\usepackage{xcolor}
\usepackage{listings}
\usepackage[portuguese]{babel}
\usepackage{microtype}
\usepackage{xspace}
\usepackage{url}
\usepackage{csquotes}
\usepackage{lipsum}
\usepackage{array}
\usepackage{multirow}
\usepackage{booktabs}
\usepackage[
    backend=bibtex,
    style=alphabetic,
    sorting=ynt,
    bibencoding=utf8
]{biblatex}
\addbibresource{Relatorio.bib}

\definecolor{keycolor}{RGB}{172, 42, 42}

\lstdefinestyle{global}{
basicstyle=\ttfamily\scriptsize\color{black!90},%
stringstyle=\itshape\color{magenta},%
showstringspaces=false,%
keywordstyle=\bfseries\color{keycolor},%
commentstyle=\color{blue}\slshape,%
framexleftmargin=1mm,%
backgroundcolor=\color{black!2},%
    }

\lstdefinestyle{makefile}{
otherkeywords={.SUFFIXES},
alsoletter={:},
morekeywords=[1]{SUFFIX, CPP_},
morekeywords=[2]{vasp:,makeparam:,zgemmtest:,dgemmtest:,ffttest:,kpoints:,clean:},
style=global,%
morecomment=[l][commentstyle]{\#},%
emphstyle={\color{vimvert}},%
moredelim=[s][\color{vimvert}]{\$(}{)}%
}

\newcommand{\head}[1]{\textnormal{\textbf{#1}}}
\newcolumntype{P}[1]{>{\centering\arraybackslash}p{#1}}

\definecolor{dkgreen}{rgb}{0,0.6,0}
\definecolor{mauve}{rgb}{0.58,0,0.82}
\definecolor{gray}{rgb}{0.5,0.5,0.5}
\lstset{basicstyle=\ttfamily,
    showstringspaces=false,
    commentstyle=\color{dkgreen},
    keywordstyle=\color{blue},
    columns=flexible,
    basicstyle={\small\ttfamily},
    numbers=left,
    numberstyle=\tiny\color{gray},
    stringstyle=\color{mauve},
    breaklines=true,
    breakatwhitespace=true,
}

% Document
\begin{document}
    \begin{figure}
        \centering
        \includegraphics[width=0.9\textwidth]{deti2}\label{fig:figure}
    \end{figure}
    \maketitle
    \tableofcontents
    \clearpage


    \section{Introdução}\label{sec:introducao}
Uma word ladder é uma sequência de palavras em que cada palavra difere em uma e só uma letra da palavra anterior.
% Very original
Por exemplo, na língua Portuguesa é possível ir da palavra tudo para a palavra nada em quatro passos.
\emph{$tudo \rightarrow todo \rightarrow nodo \rightarrow nado \rightarrow nada$}.
Como tal para resolver o problema proposto de criar um algoritmo em \emph{C/C++} que permita encontrar foi feita uma implementação de uma \emph{Hash Table} em \emph{C++}, usada depois para permitir a implementação de grafos e do \emph{union find}, tornando possível tal algoritmo.

    \section{Funções Usadas}\label{sec:funcoes-usadas}
    Esta secção contém uma lista, com a respetiva descrição de todas as funções usadas para a criação do algoritmo.

    \subsection{Hash Table}\label{subsec:hash-table}
    Sendo a linguagem de progamação escolhida para a resolução de este problema \emph{C++}, a \emph{Hash table} foi implementada através de duas classes, uma que contém os parâmetros de cada nó da \emph{Hash Table}, e uma que contém a implementação da \emph{Hash Table}.

    \subsubsection{Constructor}
    O construtor da classe hashTable é um construtor padrão que inicializa todas as variáveis da classe quando um objeto da classe é criado. Ele é chamado quando o objeto é criado sem argumentos ou com o construtor padrão.

    Ele define o tamanho da tabela hash como sendo 65536.
    Ele cria um array de ponteiros chamado "words" com o tamanho da tabela.
    Ele inicializa a variável entries para 0, a variável connected\_components para 0 e a variável load\_factor para 0,75.
    Ele define todos os elementos do array words como nullptr, o que representa que os elementos da tabela estão vazios.

    \subsubsection{Destructor}
O destrutor ele faz um loop pelo array words e deleta cada elemento não nulo e finalmente deleta o array words.
   
    \subsubsection{add}
    Esta função é usada para adicionar uma nova palavra à tabela hash. A função recebe um único argumento, uma string chamada "word" que representa a palavra a ser adicionada à tabela.

    A função começa chamando a função hash() para obter o índice em que a palavra deve ser inserida.
    Em seguida, ele verifica se a palavra já está presente na tabela verificando se o elemento no índice calculado não é nulo e se a palavra armazenada nesse elemento é a mesma da palavra a ser adicionada. Se a palavra já estiver presente, a função retorna sem adicionar a palavra à tabela.
    A função então verifica se o número de entradas na tabela mais 1 é maior ou igual ao tamanho da tabela multiplicado pelo fator de carga. Se isso for verdade, significa que a tabela está ficando cheia e a função chama a função resize() para aumentar o tamanho da tabela.
    Se o elemento no índice calculado for nulo, a função chama a função create() para criar um novo nó com a palavra e inserir no índice.
    Se o elemento no índice calculado não for nulo e a palavra armazenada nesse elemento não for a mesma da palavra a ser adicionada, a função entra em um loop while. O loop while continua até encontrar um espaço vazio na tabela ou a palavra já estiver na tabela. Para encontrar o próximo espaço disponível, ele usa o sondagem linear. A sondagem linear é uma técnica de resolução de colisão onde o próximo espaço é encontrado incrementando o índice em 1 e tomando o resto quando dividido pelo tamanho da tabela.
    Se um espaço vazio for encontrado, a função chama a função create() para criar um novo nó com a palavra e inserir no índice.

Observe que essa implementação usa o endereçamento aberto como técnica de resolução de colisão, mais especificamente a sondagem linear.

    \subsubsection{get}
    Esta função é utilizada para obter um nó específico dado uma palavra. Ela começa chamando a função hash para obter o índice onde a palavra deve ser encontrada. Se o elemento no índice calculado é nulo, a função retorna nullptr, indicando que a palavra não está presente. Se o elemento não é nulo, ele verifica se a palavra armazenada é a mesma da palavra que se deseja procurar. Se for, a função retorna o endereço desse elemento. Se não é a mesma, a função entra em um loop while e usa o sondagem linear para encontrar a palavra ou um espaço vazio. Se a palavra é encontrada, a função retorna o endereço desse elemento. Se não é encontrada, a função retorna nullptr.

    \subsubsection{add\_edge}
    Esta função é utilizada para adicionar uma aresta entre dois nós na estrutura de dados que representa um grafo. Ela recebe dois argumentos, "from" e "to", que são ponteiros para os nós entre os quais a aresta será adicionada.

    A função adiciona o nó "to" à lista de adjacência do nó "from" e o nó "from" à lista de adjacência do nó "to", estabelecendo assim a conexão entre eles.
    A função incrementa o contador de arestas em ambos os nós, indicando que eles têm uma aresta adicional.
    A função chama outra função chamada "g\_union" passando "from" e "to" como argumentos, essa função fará algo relacionado com a estrutura de dados union-find.

    \subsubsection{BFS}
    A função BFS é uma implementação da busca em largura. Ela começa percorrendo todos os nós e marcando-os como não visitados e sem pais. Em seguida, a função adiciona o nó "from" a uma fila e marca-o como visitado. A função então entra em um loop enquanto a fila não estiver vazia. Dentro do loop, a função pega o primeiro elemento da fila e verifica se é o nó "to". Se for, a função retorna a profundidade atual. Caso contrário, a função percorre todos os nós adjacentes ao nó atual que ainda não foram visitados, marca-os como visitados e adiciona-os à fila. A profundidade é incrementada a cada iteração do loop. Se a profundidade atual é maior que o valor máximo de profundidade especificado e o valor máximo de profundidade é diferente de 0, a função retorna -1. Se o loop termina e o nó "to" ainda não foi encontrado, a função retorna -1.

    \subsubsection{DFS}
    A função DFS é uma implementação da busca em profundidade. Ela funciona de maneira semelhante à função BFS, mas usa uma pilha em vez de uma fila. A função marca todos os nós como não visitados e sem pais e adiciona o nó "from" à pilha. A função então entra em um loop enquanto a pilha não estiver vazia. Dentro do loop, a função pega o topo da pilha e verifica se é o nó "to". Se for, a função segue o caminho de volta ao nó "from" contando a profundidade e retorna-a. Caso contrário, a função percorre todos os nós adjacentes ao nó atual que ainda não foram visitados, marca-os como visitados e adiciona-os à pilha. Se o loop termina e o nó "to" ainda não foi encontrado, a função retorna -1.
    
    \subsubsection{list\_connected\_components}
    A função list\_connected\_components é usada para listar todos os nós que fazem parte do mesmo componente conectado de um nó específico dado uma palavra. Ela chama a função get() para obter o nó correspondente à palavra fornecida, se ele não for encontrado, a função imprime "Palavra não encontrada" e retorna. Caso contrário, ela chama outra função chamada "find", passando o nó encontrado como argumento, essa função retorna um representante de um conjunto na estrutura de dados union-find. Em seguida, a função percorre todos os nós na tabela hash e adiciona a um vetor de componentes todos os nós cujo representante é o mesmo do nó encontrado anteriormente. Finalmente, a função imprime "Pertencente ao mesmo componente conectado como [palavra fornecida]:" e imprime todas as palavras armazenadas nos nós do vetor de componentes.
    
    \subsubsection{find}
    
    \subsubsection{g\_union}
    
    \subsubsection{print\_adjacency\_list}
    
    \subsubsection{hash}
    
    \subsubsection{unhash}


    \subsection{Estatísticas da Hash Table}\label{subsec:estatisticas-da-hash-table}

    \subsubsection{get\_load\_factor e get\_collisions}
   As funções get\_load\_factor e get\_collisions são usadas para obter estatísticas sobre a tabela hash. A primeira função retorna o fator de carga atual, que é o número de entradas na tabela dividido pelo tamanho da tabela. A segunda função retorna o número de colisões na tabela, que é o número de entradas que estão em uma posição diferente daquela calculada pela função de hash.
    
    \subsubsection{get\_distribution}
    A função get\_distribution é usada para obter a distribuição de entradas na tabela hash. Ela percorre toda a tabela e adiciona "true" a um vetor se uma entrada está presente na posição atual ou "false" se não estiver. A função retorna esse vetor.
    
    \subsection{Estatísticas do Grafo}\label{subsec:estatisticas-do-grafo}
    
    \subsubsection{get\_connected\_components}
    A função get\_connected\_components conta o número de componentes conectados no grafo. Ele percorre a tabela hash e verifica, para cada nó, se é um representante de seu conjunto de união (usando a propriedade representative). Se for, incrementa o contador de componentes. O número de componentes conectados é retornado no final.
    
    \subsubsection{get\_diameter}
    A função get\_diameter calcula o diâmetro do grafo. O diâmetro é o caminho mais longo entre dois nós no grafo. Começa por percorrer a hash table e chama a função DFS em cada nó, passando o nó atual como início e o nó especificado como parâmetro como destino. A distância retornada pela DFS é comparada com a distância máxima encontrada até agora. Se for maior, atualiza a distância máxima e salva o nó inicial. No final, se print for verdadeiro, imprime o diâmetro e o caminho encontrado.
    
    \subsubsection{get\_diameter\_node}
    A função get\_diameter\_node é semelhante à anterior, mas apenas retorna o nó inicial do caminho com o maior diâmetro.
    

    \subsection{Outras Funções}\label{subsec:outras-funcoes}

    \subsubsection{longest}
    A função "longest" encontra o caminho mais longo entre uma palavra dada e todas as outras palavras na tabela de hash. Ela faz isso chamando o método get\_diameter() da tabela de hash correspondente ao tamanho da palavra dada.
    
    \subsubsection{connected}
    A função "connected" verifica se as duas palavras dadas são conectadas (se elas diferem apenas por uma letra). Ela faz isso comparando cada letra das palavras uma a uma e verificando se elas são diferentes.
    
    \subsubsection{path\_finder}
    A função "path\_finder" é usada para encontrar o caminho mais curto entre duas palavras no grafo. Ele usa a função BFS para encontrar a distância mais curta entre as duas palavras e imprime o caminho encontrado.
    
    \subsubsection{connected\_components}
A função "connected\_components" encontra todas as palavras que estão conectadas (que diferem apenas por uma letra) à palavra dada. Ela faz isso chamando o método list\_connected\_components() da tabela de hash correspondente ao tamanho da palavra dada.
    
    \section{Resultados}\label{sec:resultados}

    \clearpage


    \section{Referências}\label{sec:referencias}~\nocite{aed,cpp}
    \printbibliography[
        heading=none,
    ]

    \clearpage


    \section{Apêndice}\label{sec:apendice}
    
    \subsection{word\_ladder.cpp}    
    \begin{lstlisting}[language=c++,label={lst:lstlisting}]
#include <algorithm>
#include <fstream>
#include <iostream>
#include <string>
#include <cmath>
#include <vector>
#include <queue>
#include <stack>
#include <thread>

using namespace std;
#define _max_word_size_ 32

class node {
public:
    node(const string &word) : word(word) {
        parent = nullptr;
        visited = false;
        representative = this;
        vertices = 1;
        edges = 0;
    }

    string word;
    // search relevant data
    node *parent;
    bool visited;
    // graph data structure
    vector<node *> adjacency_list;
    // union-find data structure
    node *representative;
    int vertices;
    int edges;
};

class hashTable {
public:
    unsigned int size;
    node **words;
    unsigned int entries;
    int connected_components;
    double load_factor;

    hashTable() {
        // Makes the dict only need to be resized once.
        size = 65536;
        words = new node *[size];
        entries = 0;
        connected_components = 0;
        load_factor = 0.75;
        for (unsigned int i = 0; i < size; i++) {
            words[i] = nullptr;
        }
    }

    ~hashTable() {
        for (unsigned int i = 0; i < size; i++) {
            if (words[i] != nullptr) {
                delete words[i];
            }
        }
        delete[] words;
    }

    void add(const string &word) {
        unsigned int index = hash(word);
        if (words[index] != nullptr && words[index]->word == word) {
            return;
        }
        if (entries + 1 >= size * load_factor) {
            resize();
        }
        if (words[index] == nullptr) {
            create(index, word);
        } else if (words[index]->word != word) {
            while (words[index] != nullptr && words[index]->word != word) {
                // Linear probing is the fastest way.
                // Probably because it uses the cache more efficiently.
                // And that matters the most when the table is huge and we have memory to spare.
                index = (index + 1) % size;
            }
            create(index, word);
        }
    }

    node *get(const string &word) {
        unsigned int index = hash(word);
        if (words[index] == nullptr) {
            return nullptr;
        }
        if (words[index]->word == word) {
            return words[index];
        }
        while (words[index] != nullptr && words[index]->word != word) {
            index = (index + 1) % size;
        }
        return words[index];
    }

    // graph functions
    void add_edge(node *from, node *to) {
        from->adjacency_list.push_back(to);
        to->adjacency_list.push_back(from);
        from->edges++;
        to->edges++;
        g_union(from, to);
    }

    int BFS(node *from, node *to, int maximum_depth = 0) {
        for (unsigned int i = 0; i < size; i++) {
            if (words[i] != nullptr) {
                words[i]->visited = false;
                words[i]->parent = nullptr;
            }
        }
        queue < node * > q;
        from->visited = true;
        q.push(from);
        int depth = 0;
        while (!q.empty()) {
            int q_size = q.size();
            for (int i = 0; i < q_size; i++) {
                node *current = q.front();
                q.pop();
                if (current == to) {
                    return depth;
                }
                for (size_t j = 0; j < current->adjacency_list.size(); j++) {
                    node *adjacent = current->adjacency_list[j];
                    if (!adjacent->visited) {
                        adjacent->visited = true;
                        adjacent->parent = current;
                        q.push(adjacent);
                    }
                }
            }
            depth++;
            if (depth > maximum_depth && maximum_depth != 0) {
                return -1;
            }
        }
        return -1;
    }

    int DFS(node *from, node *to) {
        for (unsigned int i = 0; i < size; i++) {
            if (words[i] != nullptr) {
                words[i]->visited = false;
                words[i]->parent = nullptr;
            }
        }
        stack < node * > q;
        from->visited = true;
        q.push(from);
        int depth = 0;
        while (!q.empty()) {
            int q_size = q.size();
            for (int i = 0; i < q_size; i++) {
                node *current = q.top();
                q.pop();
                if (current == to) {
                    // god why
                    while (current->parent != nullptr && current != from) {
                        current = current->parent;
                        depth++;
                    }
                    return depth;
                }
                for (size_t j = 0; j < current->adjacency_list.size(); j++) {
                    node *adjacent = current->adjacency_list[j];
                    if (!adjacent->visited) {
                        adjacent->visited = true;
                        adjacent->parent = current;
                        q.push(adjacent);
                    }
                }
            }
        }
        return -1;
    }

    void list_connected_components(const string &word) {
        vector < node * > components;
        node *vertex = get(word);
        if (vertex == nullptr) {
            cout << "Word not found" << endl;
            return;
        }
        node *representative = find(vertex);
        for (unsigned int i = 0; i < size; i++) {
            if (words[i] != nullptr && find(words[i]) == representative) {
                components.push_back(words[i]);
            }
        }
        cout << "Belonging to same connected component as " << word << "are:" << endl;
        for (size_t i = 0; i < components.size(); i++) {
            cout << components[i]->word << "\n";
        }
    }

    // hash table statistics
    double get_load_factor() {
        return (double) entries / size;
    }

    int get_collisions() {
        unsigned int collisions = 0;
        for (unsigned int i = 0; i < size; i++) {
            if (words[i] != nullptr) {
                if (hash(words[i]->word) != i) {
                    collisions++;
                }
            }
        }
        return collisions;
    }

    vector<bool> get_distribution() {
        vector<bool> distribution;
        for (unsigned int i = 0; i < size; i++) {
            if (words[i] != nullptr) {
                distribution.push_back(true);
            } else {
                distribution.push_back(false);
            }
        }
        return distribution;
    }

    // graph statistics
    int get_connected_components() {
        int components = 0;
        for (unsigned int i = 0; i < size; i++) {
            if (words[i] != nullptr) {
                if (words[i]->representative == words[i]) {
                    components++;
                }
            }
        }
        return components;
    }

    int get_diameter(node *n, bool print = true) {
        int diameter = 0;
        node *max = nullptr;
        for (unsigned int i = 0; i < size; i++) {
            if (words[i] != nullptr) {
                if (words[i]->adjacency_list.size() == 0) {
                    continue;
                }
                int distance = DFS(words[i], n);
                if (distance > diameter) {
                    diameter = distance;
                    max = words[i];
                }
            }
        }
        // DFS data is wiped out every run.
        DFS(n, max);
        node *res = max;
        if (res == nullptr) {
            return 0;
        }
        if (print) {
            cout << "Diameter: " << diameter << endl;
            cout << "Path: ";
            if (res == nullptr) {
                cout << "No connected words." << endl;
            }
            while (res->parent != nullptr) {
                cout << res->word << " -> ";
                res = res->parent;
            }
            cout << res->word << endl;
        }

        return diameter;
    }

    node *get_diameter_node(node *n) {
        int diameter = 0;
        node *max = nullptr;
        for (unsigned int i = 0; i < size; i++) {
            if (words[i] != nullptr) {
                int distance = DFS(words[i], n);
                if (distance > diameter) {
                    diameter = distance;
                    max = words[i];
                }
            }
        }
        return max;
    }

private:
    void create(int index, const string &word) {
        entries++;
        connected_components++;
        words[index] = new node(word);
    }

    void resize() {
        // High resize coefficient to reduce resizes, which are expensive.
        int coeff = 4;
        size *= coeff;
        node **new_words = new node *[size];
        for (unsigned int i = 0; i < size; i++) {
            new_words[i] = nullptr;
        }
        for (unsigned int i = 0; i < size / coeff; i++) {
            if (words[i] != nullptr) {
                int index = hash(words[i]->word);
                if (new_words[index] == nullptr) {
                    new_words[index] = words[i];
                } else {
                    while (new_words[index] != nullptr) {
                        index = (index + 1) % size;
                    }
                }
            }
        }
        delete[] words;
        words = new_words;
    }

    node *find(node *vertex) {
        if (vertex->representative != vertex) {
            vertex->representative = find(vertex->representative);
        }
        return vertex->representative;
    }

    void g_union(node *from, node *to) {
        node *from_rep = find(from);
        node *to_rep = find(to);
        if (from_rep != to_rep) {
            to_rep->representative = from_rep;
            connected_components--;
        }
    }

    void print_adjacency_list(node *n) {
        cout << n->word << " -> ";
        for (size_t i = 0; i < n->adjacency_list.size(); i++) {
            cout << n->adjacency_list[i]->word << " ";
        }
        cout << endl;
    }

#define FNV_OFFSET 14695981039346656037UL
#define FNV_PRIME 1099511628211UL

    // Return 64-bit FNV-1a hash for key (NUL-terminated). 
    unsigned int hash(const string &word) {
        uint64_t hash = FNV_OFFSET;
        const char *key = word.c_str();
        for (const char *p = key; *p; p++) {
            hash ^= (uint64_t)(unsigned char)(*p);
            hash *= FNV_PRIME;
        }
        // Ensure hash is adjusted to the size of the table.
        return (size_t)(hash & (uint64_t)(size - 1));
    }

    // https://github.com/skeeto/hash-prospector#three-round-functions
    // Kept for reference.
    unsigned int hash(int x) {
        x ^= x >> 17;
        x *= 0xed5ad4bb;
        x ^= x >> 11;
        x *= 0xac4c1b51;
        x ^= x >> 15;
        x *= 0x31848bab;
        x ^= x >> 14;
        return x;
    }

    unsigned int unhash(int x) {
        x ^= x >> 14 ^ x >> 28;
        x *= 0x32b21703;
        x ^= x >> 15 ^ x >> 30;
        x *= 0x469e0db1;
        x ^= x >> 11 ^ x >> 22;
        x *= 0x79a85073;
        x ^= x >> 17;
        return x;
    }
};

void longest(hashTable **dicts, const string &word) {
    hashTable *dict = dicts[word.size() - 1];
    node *n = dict->get(word);
    cout << "Longest path to " << word << " is " << endl
         << dict->get_diameter(n)
         << " words long." << endl;
    return;
}

bool connected(const string &a, const string &b) {
    if (a.size() != b.size())
        return false;
    bool result = false;
    for (size_t i = 0; i < a.size(); i++) {
        if (a[i] != b[i]) {
            // Only one difference is allowed
            if (result)
                return false;
            result = true;
        }
    }
    return result;
}

void path_finder(hashTable **dicts, const string &start, const string &end) {
    if (start.size() != end.size()) {
        cout << "Cannot compare different sizes." << endl;
        return;
    }
    hashTable *dict = dicts[start.size() - 1];
    cout << "Trying to go from " << start << " to " << end << endl;
    node *from = dict->get(end);
    node *to = dict->get(start);
    if (from == nullptr || to == nullptr) {
        cout << "No path found." << endl;
        return;
    }
    int travelled = dict->BFS(from, to);
    cout << "Travelled " << travelled << " nodes. " << endl;
    node *res = to;
    while (res->parent != nullptr) {
        cout << res->word << " -> ";
        res = res->parent;
    }
    cout << res->word << endl;
}

void connected_components(hashTable **dicts, const string &word) {
    hashTable *dict = dicts[word.size() - 1];
    dict->list_connected_components(word);
}

void end(hashTable **dicts) {
#if defined(_stats_) || defined(_detail_) || defined(_full_)
    ofstream file;
    file.open("stats.txt");
#endif
    for (size_t i = 0; i < _max_word_size_; i++) {
#if defined(_stats_) || defined(_detail_) || defined(_full_)
        file << endl;
        file << "Hash Table for " << i + 1 << " letter words" << endl;
        file << "Size: " << dicts[i]->size << endl;
        file << "Load factor: " << dicts[i]->get_load_factor() << endl;
        file << "Collisions: " << dicts[i]->get_collisions() << endl;
#if defined(_detail_) || defined(_full_)
        vector<bool> distribution = dicts[i]->get_distribution();
        file << "Distribution: " << endl;
        for (size_t j = 0; j < distribution.size(); j++)
        {
            if (distribution[j])
                file << j << " ";
        }
        file << endl;
#endif
#endif
        delete dicts[i];
    }
#if defined(_stats_) || defined(_detail_) || defined(_full_)
    file.close();
#endif
}

void graph_builder(hashTable *dict) {
    int sizes = 0;
    // TODO: Optimize this, O(n^1.5) ish isn't good
    for (size_t i = 0; i < dict->size; i++) {
        node *from = dict->words[i];
        if (from == nullptr)
            continue;
        if (sizes == 0)
            sizes = from->word.size();
        for (size_t j = i + 1; j < dict->size; j++) {
            node *to = dict->words[j];
            if (to == nullptr)
                continue;
            if (connected(from->word, to->word)) {
                dict->add_edge(from, to);
            }
        }
    }
    if (sizes != 0)
        cout << "Processed " << sizes + 1 << " letter words" << endl;
}

void longest_path(hashTable *dict) {
    int largest = 0;
    vector < node * > reprs;
    node *max = nullptr;
    for (unsigned int i = 0; i < dict->size; i++) {
        if (dict->words[i] != nullptr) {
            if (find(reprs.begin(), reprs.end(), dict->words[i]->representative) == reprs.end()) {
                reprs.push_back(dict->words[i]->representative);
                int depth = dict->get_diameter(dict->words[i], false);
                if (depth > largest) {
                    largest = depth;
                    max = dict->words[i];
                }
            }
        }
    }
    node *origin = dict->get_diameter_node(max);
    if (origin == nullptr || max == nullptr) {
        cout << "No path found." << endl;
        return;
    }
    dict->DFS(origin, max);
    node *res = max;
    ofstream file;
    file.open("longest.txt", ios::app);
    file << "Longest path for " << max->word.size() << " letter words" << endl;
    file << "Size: " << largest << endl;
    while (res->parent != nullptr) {
        file << res->word << " -> ";
        res = res->parent;
    }
    file << res->word << endl;
}

int main() {
    setlocale(LC_ALL, ".UTF8");
    hashTable *dicts[_max_word_size_];
    thread threads[_max_word_size_];
    for (size_t i = 0; i < _max_word_size_; i++) {
        dicts[i] = new hashTable;
    }
    ifstream in("wordlist-big-latest.txt");
    if (!in) {
        printf("Error: could not open words file\n");
    }
    string word;
    while (in >> word) {
        int size = word.size();
        dicts[size - 1]->add(word);
    }
    for (int sizes = 0; sizes < _max_word_size_; sizes++) {
        hashTable *dict = dicts[sizes];
        threads[sizes] = thread(graph_builder, dict);
    }
    for (int sizes = 0; sizes < _max_word_size_; sizes++) {
        threads[sizes].join();
    }
    path_finder(dicts, "etano", "sitie");
#ifdef _full_
    ofstream file;
    file.open("longest.txt", ios::trunc);
    file.close();
    for (int sizes = 0; sizes < _max_word_size_; sizes++)
    {
        hashTable *dict = dicts[sizes];
        threads[sizes] = thread(longest_path, dict);
    }
    for (int sizes = 0; sizes < _max_word_size_; sizes++)
    {
        threads[sizes].join();
    }
#endif
    // connected_components(dicts, "belo");
    longest(dicts, "etano");

    // TODO: See graphs
    // TODO: Interesting diameters
    // etano and sitia are opposite extremeties of (one of the) main connected component, as they show up in lots of diameters
    end(dicts);
}
    \end{lstlisting}
    
    \subsection{makefile}
        \begin{lstlisting}[style=makefile,label={lst:lstlisting}]
#
# makefile to compile the A.02 assignment (word ladder)
#

clean:
	rm -rf a.out word_ladder *.exe

word_ladder: word_ladder.cpp
	g++ -Wall -Wextra -O3 word_ladder.cpp -o word_ladder -lm -march=native

stats: word_ladder.cpp
	g++ -Wall -Wextra -O3 word_ladder.cpp -o word_ladder -lm -march=native -D_stats_

detail: word_ladder.cpp
	g++ -Wall -Wextra -O3 word_ladder.cpp -o word_ladder -lm -march=native -D_detail_

full: word_ladder.cpp
	g++ -Wall -Wextra -O3 word_ladder.cpp -o word_ladder -lm -march=native -D_full_

debug: word_ladder.cpp
	g++ -Wall -Wextra -O0 -ggdb3 word_ladder.cpp -o word_ladder -lm -march=native -D_full_
\end{lstlisting}



\end{document}
